\chapter{Hash Functions}
\label{short:hash-functions}

Can we express something about a secret without revealing it?

A digest is a small, repeatable fingerprint of exact bytes. The same bytes
produce the same digest. This offers a useful instrument for larger
constructions, but it does not answer the opening question by itself.

If the secret is one of three moves, an observer can hash all three and
compare. The secret was hidden from sight, yet remained available to a
guess.

Attend to the whole construction: the possible inputs, their encoding, the
reference value, and what an adversary can try. A compact expression is not
necessarily a private one.

\shortlimit{A bare digest neither hides a guessable input nor authenticates its sender.}
\companionref{A}{appendix.A}{Definitions, attack examples, and the BLAKE2b C demonstration.}
